\documentclass[11pt]{article}
\usepackage[margin=1in]{geometry}
\usepackage{amsmath,amssymb}
\usepackage[T1]{fontenc}
\usepackage{lmodern}
\usepackage{microtype}
\usepackage{xurl}
\usepackage{booktabs}
\usepackage{graphicx}
\usepackage{array}
\usepackage{hyperref}
\setlength{\emergencystretch}{3em}

\title{Utility, Model Differences, and Code Experiment\\Fertility and Housing Supply}
\author{Fertility and Housing Supply Project}
\date{February 20, 2026}

\begin{document}
\maketitle

\section*{Purpose}
This document gives a detailed model explanation focused on what changes relative to the baseline NIMBY paper, why those changes matter, and what was learned from a first coding experiment.

\section{Baseline NIMBY model block (reference point)}
The original baseline (as used in the NIMBY project) can be summarized as:
\begin{enumerate}
\item Household utility over consumption and housing services (CRRA/Cobb-Douglas).
\item Life-cycle decision problem for tenure, assets, housing, and consumption.
\item Housing market/political equilibrium where supply outcomes (prices/quantities) are shaped by voting incentives.
\item Exogenous demography in transition calculations.
\end{enumerate}

The fertility project keeps (1)--(3) as much as possible, but changes (4) and links household fertility choices to housing and political dynamics.

\section{Exact differences from the NIMBY paper}
\subsection*{Difference A: Utility now includes children-at-home and crowding}
\textbf{Baseline flow utility:}
\[
u(c_t,h_t)=\frac{(c_t^{1-\zeta}h_t^{\zeta})^{1-\gamma}}{1-\gamma}.
\]

\textbf{Fertility extension flow utility:}
\[
 u_t = \frac{(c_t^{1-\zeta}h_{t,eff}^{\zeta})^{1-\gamma}}{1-\gamma} + \chi\frac{n_{t,home}^{1-\eta}}{1-\eta},
\quad
h_{t,eff}=\frac{h_t}{(1+\lambda n_{t,home})^{\psi}}.
\]

\textbf{Economic meaning:}
\begin{itemize}
\item Housing enters utility through \emph{effective} space, so crowding makes the same dwelling less valuable for larger families.
\item Fertility is not only a budget burden; it can generate direct utility via the child term.
\item This creates a quantity--space tradeoff even before adding child quality spending.
\end{itemize}

\subsection*{Difference B: Fertility is now a household choice variable}
In fertile ages, births are chosen jointly with consumption/housing/saving:
\[
V_t(s_t)=\max_{c_t,a_{t+1},h_{t+1},b_t}\{u_t+\beta E_tV_{t+1}(s_{t+1})\}
\]
subject to the extended budget:
\[
 c_t + q_t^h h_{t+1} + a_{t+1} + \kappa(b_t,n_{t,home})
\le y_t + (1+r_t)a_t + \text{housing cash-flow terms}.
\]

\textbf{Economic meaning:}
\begin{itemize}
\item In the baseline, demographics do not respond to housing affordability.
\item In the extension, higher housing costs can delay or reduce births through both utility crowding and budget channels.
\end{itemize}

\subsection*{Difference C: Children leave home explicitly}
Children-at-home state evolves with cohort aging and a leave-home rule (deterministic benchmark):
\begin{align*}
N_{0,t+1} &= b_t, \\
N_{a+1,t+1} &= N_{a,t}, \\
n_{t,home} &= \sum_{a=0}^{A_{leave}-1}N_{a,t}.
\end{align*}

\textbf{Economic meaning:}
\begin{itemize}
\item Utility and housing demand respond to \emph{dependent} children, not lifetime births.
\item This allows realistic life-cycle housing demand reversals (space pressure rises then later falls).
\end{itemize}

\subsection*{Difference D: Demography becomes endogenous}
Cohort masses evolve with endogenous births:
\begin{align*}
M_{j+1,t+1} &= \ell_j M_{j,t}, \\
M_{1,t+1} &= \sum_{j\in\mathcal F} f_j M_{j,t}.
\end{align*}

\textbf{Economic meaning:}
\begin{itemize}
\item Fertility decisions alter future age distribution.
\item Age distribution alters future median-voter composition.
\item Median-voter composition alters future supply outcomes.
\end{itemize}

\subsection*{Difference E: Equilibrium is now a joint price--demography fixed point}
Instead of solving prices/policies with exogenous demographic weights, the extension solves:
\[
\{p_t,r_t,\text{policy}_t,M_t\}_{t\ge 0}
\]
jointly, because each object feeds back into the others over time.

\section{Why this matters for your paper contribution}
The key conceptual gain is that the model now has a full mechanism from housing supply restrictions to fertility outcomes:
\[
\text{planning/supply} \rightarrow \text{housing prices + unit mix} \rightarrow \text{crowding + fertility choice} \rightarrow \text{future demography}.
\]
This is stronger than treating fertility as an exogenous reduced-form trend because it ties fertility outcomes to policy-relevant housing margins.

\section{Code experiment: objective and design}
\subsection*{Objective}
The prototype was not intended for calibration or quantitative policy claims. Its purpose was to test numerical behavior of the new fixed-point structure and confirm whether the convergence concerns in the model setup are operationally important.

\subsection*{Prototype architecture}
A compact simulation was implemented with:
\begin{itemize}
\item age shares (8 cohorts),
\item endogenous births in fertile cohorts,
\item a political-tightness index increasing with older-share and decreasing with younger-share,
\item price updates from demand--supply imbalance,
\item optional family-sized supply shock,
\item outer-loop damping parameter.
\end{itemize}

\subsection*{Scenarios run}
\begin{enumerate}
\item No damping (aggressive update).
\item Damping $\xi=0.25$ (stabilized update).
\item Damping $\xi=0.25$ plus family-sized supply shift from period 15.
\end{enumerate}

\section{Code experiment results and diagnostics}
\subsection*{Main diagnostics}
\begin{center}
\begin{tabular}{lccc}
\toprule
Scenario & Avg price (t=10..49) & Avg fertility (t=10..49) & Max $|\Delta p_t|$ \\
\midrule
No damping & -5.939 & 0.723 & 2.426 \\
Damping 0.25 & -2.597 & 0.416 & 1.948 \\
Damping 0.25 + policy & -2.638 & 0.443 & 1.349 \\
\bottomrule
\end{tabular}
\end{center}

\textbf{Interpretation:}
\begin{itemize}
\item Undamped updates generate large jumps and near-corner behavior.
\item Damping reduces oscillation and keeps trajectories in a more interpretable range.
\item In the damped setting, family-sized supply support increases fertility and lowers transition volatility.
\end{itemize}

\begin{figure}[h!]
\centering
\includegraphics[width=0.9\textwidth]{build/model_experiment_paths.png}
\caption{Prototype dynamics under alternative update rules and policy shift}
\end{figure}

\subsection*{Failure log (important)}
Early runs produced exponential overflow and pathological paths. Two stabilization changes were required:
\begin{enumerate}
\item Bounding exponent arguments in demand/supply transforms.
\item Clamping price updates and using damped outer-loop updates.
\end{enumerate}
This is direct evidence that the full MATLAB implementation should start with tight numerical guards.

\section{Implementation implications for the full model}
\subsection*{What should be coded first}
\begin{enumerate}
\item Household block with fertility control and children-at-home state.
\item Demographic transition block that consumes policy functions and updates cohort masses.
\item Outer fixed-point wrapper with damping and share normalization.
\end{enumerate}

\subsection*{Recommended convergence checks}
Use three tolerances, not one:
\begin{itemize}
\item price-path distance,
\item cohort-share distance,
\item aggregate births distance.
\end{itemize}
All three should converge before declaring equilibrium solved.

\subsection*{Suggested baseline robustness grid}
\begin{itemize}
\item Damping values: $\xi\in\{0.10,0.25,0.40\}$.
\item Leave-home age: $A_{leave}\in\{18,19,21\}$.
\item Child-cost curvature and crowding intensity: low/medium/high triplets.
\end{itemize}

\section*{Bottom line}
The fertility extension is not a marginal tweak to the NIMBY model; it changes the equilibrium object from a housing-policy fixed point with exogenous demographics to a joint housing--demography fixed point. The prototype confirms that this is computationally feasible but numerically delicate, so disciplined damping and normalization should be treated as core model ingredients.

\end{document}
